\documentclass[11pt]{article}
\usepackage[margin=1in]{geometry}
\usepackage{amsmath, amssymb, amsthm, bm, mathtools}
\usepackage{booktabs}
\usepackage{array}
\usepackage{enumitem}

\title{Tutorial 2}
\author{ECON 101 -- Macroeconomics}
\date{Fall 2025}

\begin{document}
\maketitle

\section*{Instructions}
Answer the following questions carefully. Show all calculations and provide brief explanations where required. Use correct units (billions, millions, etc.). 

\bigskip

\section*{Problem 1: Calculating GDP with the Expenditure Approach}
Suppose the Canadian economy in 2024 has the following data (in billions of dollars):  
\begin{itemize}
    \item Consumption $(C) = \$1{,}800$  
    \item Investment $(I) = \$600$  
    \item Government purchases $(G) = \$720$  
    \item Exports $(X) = \$450$  
    \item Imports $(M) = \$520$  
\end{itemize}

\begin{enumerate}[label=\alph*)]
    \item Calculate nominal GDP using the expenditure approach.  
    \item Suppose the population is 40 million. Calculate GDP per capita.  
    \item If potential GDP is \$3{,}800 billion, compute the output gap as a percentage of potential GDP.  
\end{enumerate}

\bigskip

\section*{Problem 2: GDP from the Income Side}
You are given the following data for Canada (in billions of dollars):  
\begin{itemize}
    \item Wages: \$1{,}500  
    \item Rent: \$200  
    \item Interest: \$250  
    \item Profits: \$400  
    \item Indirect taxes (less subsidies): \$250  
    \item Depreciation: \$300  
\end{itemize}

\begin{enumerate}[label=\alph*)]
    \item Compute Net Domestic Income at factor cost.  
    \item Add indirect taxes (less subsidies) and depreciation to compute GDP at market prices.  
    \item Compare your answer with Problem 1 and explain the relationship between expenditure and income measures of GDP.  
\end{enumerate}

\bigskip

\section*{Problem 3: Nominal vs. Real GDP and Price Indexes}
The following data describe the economy’s production of two goods (apples and laptops):  

\begin{center}
\begin{tabular}{c c c c c}
\toprule
Year & Quantity of Apples & Price of Apples & Quantity of Laptops & Price of Laptops \\
\midrule
2022 & 1,000 & \$2 & 200 & \$1,000 \\
2023 & 1,200 & \$3 & 220 & \$1,100 \\
2024 & 1,300 & \$3 & 250 & \$1,200 \\
\bottomrule
\end{tabular}
\end{center}

\begin{enumerate}[label=\alph*)]
    \item Compute nominal GDP for each year.  
    \item Using 2022 as the base year, compute real GDP for 2022--2024.  
    \item Calculate the GDP deflator for each year and interpret what it measures.  
    \item What is the inflation rate between 2023 and 2024, as measured by the GDP deflator?  
\end{enumerate}

\bigskip

\section*{Problem 4: Underground and Informal Economy}
Canada’s measured GDP in 2024 is \$2{,}900 billion. Economists estimate that the underground economy accounts for 7\% of measured GDP.  

\begin{enumerate}[label=\alph*)]
    \item Compute the actual size of the underground economy in billions of dollars.  
    \item Adjust measured GDP to include the underground economy.  
    \item If the population is 39 million, what is the difference in GDP per capita between measured GDP and adjusted GDP?  
    \item Briefly explain why excluding the underground economy can distort international GDP comparisons.  
\end{enumerate}

\bigskip

\section*{Problem 5: Chain-Weighted Real GDP}
Consider the following data for a simplified economy that produces only coffee and textbooks:  

\begin{center}
\begin{tabular}{c c c c c}
\toprule
Year & Quantity Coffee & Price Coffee & Quantity Textbooks & Price Textbooks \\
\midrule
2023 & 500 & \$4 & 100 & \$60 \\
2024 & 600 & \$5 & 120 & \$70 \\
\bottomrule
\end{tabular}
\end{center}

\begin{enumerate}[label=\alph*)]
    \item Compute nominal GDP for 2023 and 2024.  
    \item Compute real GDP in 2024 using 2023 prices.  
    \item Compute real GDP in 2023 using 2024 prices.  
    \item Calculate the chain-weighted real GDP growth rate between 2023 and 2024.  
    \item Explain why the chain-weighted method is preferred over using a fixed base year.  
\end{enumerate}

\bigskip

\section*{Problem 6: GDP and Economic Well-Being}
Suppose two countries (Country A and Country B) have the same GDP per capita of \$40,000. However, the data show:  
\begin{itemize}
    \item Country A works on average 2,200 hours per worker per year.  
    \item Country B works on average 1,600 hours per worker per year.  
\end{itemize}

\begin{enumerate}[label=\alph*)]
    \item Compute GDP per hour worked for each country.  
    \item Which country has higher productivity?  
    \item Discuss why GDP per capita alone might not fully capture differences in well-being.  
\end{enumerate}

\end{document}
